\subsection{绝对值的延拓}

命题 9. 设 K 为域，L 为 K 的代数扩张，f 为 K 上的绝对值。则 f 可延拓为 L 上的绝对值。

首先假设存在一个赋值 \(v\)，它在 \(\mathbf{K}\) 上取实值，使得 \(f(x) = e^{-v(x)}\)。存在一个赋值 \(v'\)，它在 \(\mathbf{L}\) 上，其在 \(\mathbf{K}\) 上的限制与 \(v\) 等价（§ 3，第 3 节，命题 5）。于是 \(v'\) 的高为 0 或 1（第 1 节，命题 1 的推论 2），因而可设其取实值。映射 \(x \mapsto e^{-v'(x)}\) 在 \(\mathbf{K}\) 上的限制是与 \(f\) 等价的绝对值，因而具有 \(f^s\) 的形式，其中 \(s > 0\)（一般拓扑，第 IX 章，第 3 节，第 2 节，命题 5）。于是

\[
x \mapsto e ^ {- v ^ {\prime} (x) / s}
\]

是 \(\mathbf{L}\) 上延拓 \(f\) 的绝对值。

现在假设 \(f\) 不是超度量绝对值。于是 \(K\) 可识别为 \(\mathbf{C}\) 的一个子域，使得 \(f(x) = |x|^s\)，其中 \(0 \leqslant s \leqslant 1\)（§ 6，第 4 节，定理 2）。由于 \(\mathbf{C}\) 代数闭，\(L\) 可识别为 \(\mathbf{C}\) 的一个子域，且绝对值 \(x \mapsto |x|^s\) 延拓 \(f\)。

命题 10. 设 K 为域，f 为 K 上的绝对值，K 关于 f 完备且非离散，L 为 K 的代数扩张。则 f 可唯一延拓为 L 上的绝对值 \(f'\)，并且若 L 的次数为 n，则

\[
f ^ {\prime} (x) = \left(f (\mathrm{N} _ {\mathrm{L/K}} (x))\right) ^ {1 / n}
\]

对所有 \(x \in \mathbf{L}\) 成立。

\(f'\) 的存在性由命题 9 得到，而其唯一性（在 \(L\) 的每个有限子扩张上，因而在整个 \(L\) 上）由 § 6，第 4 节的引理 2 得到。令 \(f'\) 为 \(f\) 到 \(K\) 的代数闭包的唯一延拓，并

假设 L 的次数为 n。~我们知道 \(\mathrm{N}_{\mathrm{L}/\mathrm{K}}(x)=\prod_{i=1}^{n} x_i\)，其中每个 \(x_{i}\) 都是 x 在 K 上的共轭元（代数，第 VIII 章，第 12 节，第 2 节，命题 4）。由唯一性，\(f', f'(x_{i})=f'(x)\) 对所有 i 成立，从而得到所述公式。

命题 11. 设 K 为域，f 为 K 上的非超度量绝对值，\(\hat{K}\) 为 K 关于 f 的完备化，\(\hat{f}\) 为 f 到 \(\hat{K}\) 的连续延拓，L 为 K 的次数为 n 的有限扩张。

\begin{enumerate}
\def\labelenumi{(\alph{enumi})}
\item
  令 \(f'\) 为 \(\mathbf{L}\) 上延拓 \(f\) 的绝对值；令 \(\hat{\mathbf{L}}_{f'}\) 表示 \(\mathbf{L}\) 关于 \(f'\) 的完备化，并将 \(\hat{\mathbf{K}}\) 识别为 \(\mathbf{K}\) 在 \(\hat{\mathbf{L}}_{f'}\) 中的闭包；则 \([\hat{\mathbf{L}}_{f'}: \hat{\mathbf{K}}] \leqslant n\)。
\item
  L 上延拓 f 的绝对值的数目有限。若将它们记为 \(f_{1}^{\prime}, \ldots, f_{s}^{\prime}\)，并将 L 关于 \(f_{i}^{\prime}\) 的完备化记为 \(\hat{L}_{i}\)，则典范映射
\end{enumerate}

\(\hat{\mathbf{K}}\otimes_{\mathbf{K}}\mathbf{L}\rightarrow \prod_{i = 1}^{n}\hat{\mathbf{L}}_i\) 是同构，并且

\[
\sum_ {i = 1} ^ {s} \left[ \hat {\mathrm{L}} _ {i}: \hat {\mathrm{K}} \right] = n.\tag{10}
\]

证明与命题 2 中相应断言的证明相同（第 2 节）。所引用的

§ 7，第 2 节，定理 1；§ 5，第 2 节，命题 4 的推论

应替换为下列

§ 7，第 3 节，定理 2；拓扑向量空间，第一章，

第 2 节，第 3 节，定理 2 的推论 1。

注意，\(f\) 到 \(\mathbf{L}\) 的两个定义相同拓扑的延拓相等（一般拓扑，第 IX 章，第 3 节，第 2 节，命题 5）。最后，由于 \(f\) 不是超度量绝对值，\(\mathbf{K}\) 的特征为 0，因而 \(\hat{\mathbf{K}} \otimes_{\mathbf{K}} \mathbf{L}\) 的 Jacobson 根为零。

\textbf{注}\par

\begin{enumerate}
\def\labelenumi{(\arabic{enumi})}
\item
  命题 11 (b) 表明，\(\hat{\mathbf{K}}\) 与 \(\mathbf{L}\) 在 \(\mathbf{K}\) 上的每个复合扩张都同构于某个完备化 \(\hat{\mathbf{L}}_i\)，并且这些复合扩张两两不同构。
\item
  我们知道完备化 \(\hat{\mathbf{K}}\) 和 \(\hat{\mathbf{L}}_i\) 同构于 \(\mathbf{R}\) 或 \(\mathbf{C}\)（§ 6，第 4 节，定理 2）。如果 \(\hat{\mathbf{K}}\) 同构于 \(\mathbf{C}\)，那么 \(\hat{\mathbf{L}}_i\) 对所有 \(i\) 也如此，且 (10) 表明延拓 \(f_i'\) 的数目等于 \(n\)。如果 \(\hat{\mathbf{K}}\) 同构于 \(\mathbf{R}\)（例如 \(\mathbf{K} = \mathbf{Q}\)），令 \(r_1\)（分别地，\(r_2\)）表示指标 \(i\) 的数目，其中 \(\hat{\mathbf{L}}_i\) 同构于 \(\mathbf{R}\)（分别地 \(\mathbf{C}\)）；则 (10) 可写为：
\end{enumerate}

\[
r _ {1} + 2 r _ {2} = n.\tag{11}
\]

命题 12. 设 K 为域，f 为 K 上的绝对值，L 为 K 的拟 Galois 扩张，且 \(f'\) 和 \(f''\) 是 f 到 L 的两个延拓。则存在 L 的一个 K-自同构 s，使得 \(f'' = f' \circ s\)。

若 \(f\) 是超度量绝对值，命题 7 的推论 1（第 6 节）表明，存在一个 K-自同构 \(s\)，作用于 \(L\)，使得 \(f''\) 和 \(f' \circ s\) 是等价绝对值；于是存在实数 \(a > 0\)，使得 \(f''(x) = (f'(s(x)))^a\)，对所有 \(x \in L\) 成立。若 \(f\) 不是非真绝对值，取 \(x \in K^*\) 使得 \(f(x) \neq 1\)，即可知 \(a = 1\)。若 \(f\) 是非真绝对值，则 \(f'\) 和 \(f''\) 也如此（命题 1 的推论 2，第 1 节），并且 \(s\) 可取为恒等自同构。

若 \(f\) 不是超度量绝对值，则存在 \(\mathbf{Q}\)-同构 \(u', u''\)，它们将 \(\mathbf{L}\) 同构到 \(\mathbf{C}\) 的子域中，并存在实指数 \(a' > 0\)、\(a'' > 0\)，使得 \(f'(x) = |u'(x)|^{a'}\)，并且

\[
f ^ {\prime \prime} (x) = | u ^ {\prime \prime} (x) | ^ {a ^ {\prime \prime}}
\]

对所有 \(x \in \mathbf{L}\) 成立（§ 6，第 4 节，定理 2）。取 \(x = 2\)，可见 \(a' = a''\)。\(u'\) 和 \(u''\) 在 \(\mathbf{K}\) 上的限制通过连续性延拓为同构 \(u_1\) 和 \(u_2\)，从 \(\hat{\mathbf{K}}\) 到 \(\mathbf{R}\)（分别地，\(\mathbf{C}\)）。于是 \(u_2 \circ \bar{u}_1^{-1}\) 是有赋值域 \(\mathbf{R}\)（分别地，\(\mathbf{C}\)）的自同构，因而是恒等自同构（分别地，是恒等自同构或自同构 \(c: \zeta \to \bar{\zeta}\)）。必要时将 \(u'\) 替换为 \(c \circ u'\)，可设 \(u'\) 和 \(u''\) 在 \(\mathbf{K}\) 上的限制相同。借助这个公共限制，把 \(\mathbf{K}\) 识别为 \(\mathbf{C}\) 的子域；\(u'\) 和 \(u''\) 是从 \(\mathbf{L}\) 到 \(\mathbf{C}\) 的 K-同构。由于 \(\mathbf{L}\) 是 \(\mathbf{K}\) 的拟 Galois 扩张，存在一个 K-自同构 \(s\)，作用于 \(\mathbf{L}\)，使得 \(u'' = u' \circ s\)；由于 \(a' = a''\)，立即得到 \(f'' = f' \circ s\)。

注 (3)。若 \(\hat{\mathbf{K}}\) 同构于 \(\mathbf{R}\)，命题 12 表明，\(\hat{\mathbf{L}}_t\)（\(\mathbf{L}\) 的所有完备化，采用命题 11 的记号）彼此同构。因此，采用上面注 2 的记号，要么 \(r_1 = n\) 且 \(r_2 = 0\)，要么 \(r_1 = 0\) 且 \(2r_2 = n\)。

